% Generated draft: requires mathematical and source review.
\documentclass[11pt]{article}
\usepackage[margin=25mm]{geometry}
\usepackage{fontspec}
\setmainfont{Latin Modern Roman}
\usepackage{amsmath,amssymb,mathtools,unicode-math}
\setmathfont{Latin Modern Math}
\usepackage{enumitem,needspace,xurl,hyperref}
\setlength{\parindent}{0pt}
\setlength{\parskip}{5pt}
\setlength{\emergencystretch}{3em}
\newcommand{\N}{\mathbb{N}}
\newcommand{\Z}{\mathbb{Z}}
\newcommand{\Q}{\mathbb{Q}}
\newcommand{\R}{\mathbb{R}}
\newcommand{\C}{\mathbb{C}}
\newcommand{\F}{\mathbb{F}}
\newcommand{\A}{\mathbb{A}}
\newcommand{\PP}{\mathbb{P}}
\newcommand{\E}{\mathbb{E}}
\newcommand{\eps}{\varepsilon}
\newcommand{\dd}{\,\mathrm{d}}
\begin{document}
\section*{UnsolvedMath ID: 30001663}
\textbf{Source problem code:} catalogue entry
\subsection{Definitions and mathematical statement}
For this draft, use the conventional interpretation of an oriented simplex described below. The supplied record does not define ``compatible orientation'' or specify the admissible range of \(d\); these qualifications require source review.

Fix integers \(d\geq 2\) and \(n\geq d+1\), and let \(V\) be an \(n\)-element set. An orientation of a \(d\)-element set \(S\) is an ordering of its elements modulo even permutations. Write \([v_1,\ldots,v_d]\) for an orientation; an odd permutation gives its opposite, denoted \(-[v_1,\ldots,v_d]\). A \(d\)-tournament \(\tau\) on \(V\) assigns one of the two orientations to each \(S\in\binom{V}{d}\). In this interpretation these choices are independent; no global compatibility condition is imposed.

A set \(T\in\binom{V}{d+1}\) is \emph{directed} if there is an ordering \((v_0,\ldots,v_d)\) of its vertices such that
\[
\tau\bigl(T\setminus\{v_i\}\bigr)
 =(-1)^i[v_0,\ldots,\widehat{v_i},\ldots,v_d]
\qquad(0\leq i\leq d),
\]
where the hat means omission and multiplication by \(-1\) reverses orientation. Thus its assigned face orientations are the boundary orientations induced by one orientation of the simplex on \(T\).

Define
\[
D(\tau)=\#\left\{T\in\binom{V}{d+1}:T\text{ is directed}\right\},
\qquad
M_d(n)=\max_{\tau}D(\tau),
\qquad
\rho_d(n)=\frac{M_d(n)}{\binom{n}{d+1}}.
\]
The finite extremal question is to determine \(M_d(n)\) for the intended values of \(d\) and \(n\). The associated asymptotic questions are to determine, for each fixed \(d\), the limiting density denoted in the source by
\[
c_d=\lim_{n\to\infty}\rho_d(n),
\]
and to determine its behavior as \(d\to\infty\). This draft records the source's limiting-density formulation without supplying a proof of existence of that limit. Determining the limiting densities alone does not determine all finite values \(M_d(n)\), so the source's word ``equivalently'' requires this distinction.

The supplied dated assessment labels the problem partially solved and attributes \(c_2=c_3=1/4\), together with general upper and lower bounds, to the cited paper; it describes exact values for general \(d\) as still open. These are supplied literature claims, not independently verified conclusions of this draft.
\subsection{Short English statement}
Orient every \(d\)-element subset of an \(n\)-vertex set and maximize the number of \((d+1)\)-element subsets whose face orientations form an oriented simplex boundary. Determine the finite maximum, the limiting maximum proportion for fixed \(d\), and its behavior as \(d\) grows. The orientation convention and parameter range used here require confirmation against the source.
\subsection{Sources}
\begin{itemize}
\item Local UnsolvedMath registry, permanent ID \(30001663\), title ``Maximum Density of Directed Simplices in Tournaments,'' source problem code \texttt{catalogue entry}, in the Oberwolfach Reports Open Problems collection. The supplied record gives proposed year \(2011\) and does not identify a proposer.
\item Supplied report citation: \emph{Combinatorics} (2011), Oberwolfach Reports, DOI \texttt{10.4171/owr/2011/01}, pp.\ 5--83. \url{https://doi.org/10.4171/owr/2011/01}. The background describes Oberwolfach Reports as documenting mathematical workshops and their open-problem discussions. No precise problem-page location is supplied.
\item Supplied literature reference, labelled ``Leader--Tan'': \url{https://arxiv.org/abs/0906.4027}.
\item Supplied literature reference, labelled ``Mathematika publication'': \url{https://doi.org/10.1112/S0025579309000539}.
\item Literature assessment checked at \texttt{2026-08-21}: the supplied curation labels the problem partially solved, attributes \(c_2=c_3=1/4\) and general upper and lower bounds to the cited paper, and says exact \(c_d\) remains open for general \(d\). This is supplied dated curation derived from the OWR record, not independently verified here; its time-sensitive status and cited claims require independent verification.
\item The supplied registry labels the statement \texttt{corrected\_verified} and reports that the official report was checked and the statement repaired to match its passage. That is a supplied verification assertion; this draft has not independently consulted the report or literature.
\end{itemize}
\end{document}
